\section{AI in the Research Process}\label{sec:use-of-ai}

The disclosure preceding the table of contents records which parts of this
thesis involved AI and clarifies how its authorship should be understood.  This
section has a different purpose: it considers what this project suggests about
using coding agents in mathematical research.  The most consequential uses
were not isolated requests for prose or code.  They were workflows in which an
agent produced an artifact, a test or reviewer tried to expose its defects, and
the result was either revised, retained, or discarded under my direction.

This account is based on the project's Git history, recorded agent sessions,
and the surviving code, tests, and experiment artifacts.  These sources show
what was requested, executed, reviewed, and retained.  They do not record all
offline work, and Git records retained artifact states rather than the origin
of every idea.  The examples below therefore support claims about observed
workflows, not a causal estimate of how quickly the thesis would have developed
without AI.  The retrospective lessons are mine; the process record is used to
check them against what actually happened.

\subsection{Moving sooner from a specification to an artifact}

In retrospect, an important benefit was that a sufficiently explicit
mathematical or computational specification could be turned into an
inspectable artifact earlier than I had expected to implement it myself.  Early
in the project, an agent translated a public MATLAB implementation of the
algorithm of Haim-Kislev~\cite{HK2017} into Rust.  I reconstruct that within a
few days this yielded a program that ran roughly two orders of magnitude faster
than the MATLAB code and made experiments possible while I was still learning
the algorithm in detail.  It was not a trustworthy implementation: it returned
wrong results in some cases, had poor code quality, and omitted important
pruning.  Considerable later work was needed to obtain the implementation used
in this thesis.

A smaller example shows the same change in the order of work more cleanly.  In
a search involving crosspolytopes, an agent produced a running program and a
correct volume smoke test within minutes, and committed an initial version
shortly afterwards.  The long search still required redesign, backtracking,
symmetry reduction, checkpointing, and verification.  The early artifact was
valuable because it converted an abstract proposal into something that could
be tested and criticized; it was not evidence that the task was complete.

These examples suggest a conditional advantage rather than a general speedup:
when the contract is explicit and relevant infrastructure already exists,
delegation can move experimentation earlier.  The resulting artifact should be
treated as a question made executable, not as an answer certified by the cost
or speed with which it was produced.

\subsection{Production and verification are different tasks}

A recurring failure mode in the observed episodes was plausible work that
satisfied an incomplete proxy for the intended mathematical goal.  Tests and reviews helped
only when they were capable of distinguishing the relevant error from a correct
implementation.  This distinction was examined in a small controlled replay
of a historical bug in the reduced gradient of a projection routine.  Four
fresh agents received the same pre-fix code and the same description of the
mathematical discrepancy, without being given the historical fix or regression
test.  The prompts varied the model generation and whether the agent was told
explicitly to design a discriminating verifier before repairing the code.

\begin{table}[htbp]
  \centering
  \caption{Controlled replay of a sign error.  A regression was counted as
  discriminating only if an independent mutation restoring the erroneous sign
  made it fail.  ``Verifier first'' separated construction of that check from
  repair of the implementation.  This is a test of one workflow on one bug,
  not a ranking of model generations.}
  \label{tab:ai-sign-replay}
  \begin{tabular}{@{}lll@{}}
    \toprule
    Model & Prompt structure & Discriminating regression \\
    \midrule
    GPT-5.5 & minimal        & yes \\
    GPT-5.5 & verifier first & yes \\
    GPT-5.6 & minimal        & no  \\
    GPT-5.6 & verifier first & yes \\
    \bottomrule
  \end{tabular}
\end{table}

All four agents repaired the production sign.  However, the minimal GPT-5.6
run merely strengthened an existing symmetric test; that test still passed
when the bad sign was restored.  Its output therefore looked more thoroughly
tested without supplying evidence for the repair at issue.  Both verifier-first
runs produced a check that failed before and passed after the repair, as did the
minimal GPT-5.5 run.  The result does not establish that one model generation
is generally better than another.  It does show why the verification contract
remains necessary even when a newer agent can readily produce a plausible fix.

In practice, I obtained more reliable results by separating two roles.  One
role maximizes a stated goal: implement the algorithm, find an example, or
improve performance.  The other asks whether the stated goal is the right
proxy and constructs adversarial checks before accepting an apparent success.
For mathematical software, a test suite is therefore not valuable merely
because it exists.  It is evidence only for failures that its fixtures and
oracles could detect.

\subsection{Different tasks exposed different limits}

The record does not support one score for ``mathematical reasoning.''  In one
source-anchored review of a proposed flow-graph theorem, an agent found
consequential mismatches between the cited paper, the formal statement, the
Rust implementation, and the thesis claim.  The resulting corrections were
merged and narrowed the theorem's hypotheses and implementation boundary.  The
same episode did not prove the theorem, which remained explicitly unverified.
This is evidence that review found valuable defects in that case, not that
review is generally more reliable than proof generation.

The selected difficult proof-development episodes had a different outcome.  An
attempt to understand local behavior of the systolic ratio produced an
executable transition diagnostic for one failure mechanism, but not the sought
finite-neighborhood theorem.  A convergence-validity investigation showed that
computed endpoints could remain improvable, undermining their interpretation
as converged local maxima; it did not produce a replacement theorem or proof.
Thus agents helped turn informal concerns into obligations, counterchecks, and
empirical falsifications.  These selected cases do not estimate how often an
agent can prove a theorem, but they separate useful diagnostic progress from
proof completion.

Candidate generation showed a similar distinction.  A data-science proposer
exported actual polytope candidates and enriched the upper tail of the observed
systolic ratios after a control bug was corrected.  It produced neither a new
ratio above one nor a new mathematical conjecture, and the development branch
was later replaced by a maintained proposer with the same bounded status.
Generating candidates under an executable proxy was useful; deciding whether
the proxy answered the research question and whether the result deserved
integration remained a separate step.

\subsection{Breadth makes selection and integration more important}

Agents increased the number of computational or mathematical leads that could
be pursued in parallel.  A recorded data-science episode near the end of the project, for
example, branched into roughly nine object-level lines of work and thirty
structurally linked descendant sessions.  It produced thousands of lines of
candidate implementation together with no-go results, repairs, and several
preserved commits.  This work was later integrated as bounded research state,
but the principal forward experiment had not been run on its real target.
Thus output volume would badly overstate the empirical conclusion reached.

The same issue appeared at smaller scale in profiling work: technically useful
measurements were accompanied by an overconfident recommendation and a
contaminated handoff, so the factual artifact had to be separated from its
interpretation and reviewed afresh.  These episodes changed what needed scarce
attention.  As producing another implementation, search, or analysis became
easier, the important choices were increasingly which question to pursue, what
would discriminate success from a convincing failure, and whether the result
deserved the human work required for review and integration.

This does not imply that breadth is undesirable.  It supports a breadth-first
research strategy when probes are cheap, provided that stopping rules are set
and unfinished branches are not mistaken for results.  Cost of production is a
poor proxy for value in such a workflow: the relevant accounting includes my
time spent specifying, checking, correcting, and deciding what enters the
thesis.

\subsection{Lessons for future mathematical work}

The project led me to the following working principles.

\begin{enumerate}
  \item Delegate work once the mathematical contract and the observable success
  condition are clear enough to inspect.  An early artifact can be useful even
  when it is expected to be wrong.

  \item Design a discriminating check before optimizing for the desired result,
  especially for numerical code, proof search, and experiments whose apparent
  success is easy to fake accidentally.  Independent mutation, alternative
  derivations, exact special cases, and adversarial inputs were more informative
  than additional agreeable reviews.

  \item Keep a record of failures and of what the contemporary review did or
  did not catch.  When a session's reasoning or worktree becomes contaminated,
  restarting is often cheaper than repairing its trajectory; independently
  valid artifacts and a short account of the failure can still be carried
  forward.

  \item Evaluate a portfolio by retained evidence and human attention, not by
  generated code, session count, or apparent labor saved.  Increased breadth is
  valuable only when weak lines are stopped and promising ones are integrated.

  \item Re-test workflow assumptions after a model change.  Tactics for reducing
  friction may become obsolete quickly, whereas the need to separate production
  from discriminating verification survived the model-generation comparison
  performed here.
\end{enumerate}

AI therefore affected this project less by supplying a single mathematical
idea than by changing the feasible research loop: more proposals could be made
executable, criticized, and abandoned.  That loop was productive when the
verification target and selection decision remained separate from the agent's
task of producing a persuasive result.  The responsibility for those targets,
and for every claim retained in the thesis, remained mine.
